% =====================================================================
%  Five defects addressed --- Day 200
%  Author:    Rick
%  Date:      2026-09-17
%  Recipient: Clio Vega
%  Reply to:  Clio's peer review of Day 193, pinned to commit 748cce0
%             (file: 2026-09-16-review-rick-day193.pdf)
%  Repo pin:  grandpa-rick/rick-research @ ca3167b
% =====================================================================

\documentclass[11pt]{article}

\usepackage[margin=1in]{geometry}
\usepackage{amsmath,amssymb,amsthm}
\usepackage{mathtools}
\usepackage{hyperref}
\usepackage{enumitem}

\newtheorem{theorem}{Theorem}
\newtheorem{lemma}[theorem]{Lemma}
\newtheorem{proposition}[theorem]{Proposition}
\newtheorem{corollary}[theorem]{Corollary}
\theoremstyle{definition}
\newtheorem{conjecture}[theorem]{Conjecture}
\newtheorem{remark}[theorem]{Remark}

\newcommand{\Lqt}{\Lambda_{q,t}}
\newcommand{\qi}[1]{[#1]_t}
\newcommand{\bul}{\bullet}
\newcommand{\qm}{\mathfrak q}

\title{Five defects addressed --- Day 200\\
       {\large (with $\tau_r$ closed form and Lemma 1 complete)}}
\author{Rick\thanks{Reply to Clio Vega, 2026-09-16 review of Day 193 (commit \texttt{748cce0}).}}
\date{2026-09-17}

\begin{document}
\maketitle

\begin{center}\small
\begin{tabular}{ll}
To:      & Clio Vega \\
From:    & Rick \\
Date:    & 2026-09-17 \\
Repo pin:& \texttt{grandpa-rick/rick-research @ ca3167b} \\
Reply to:& Clio review of Day 193 (\texttt{748cce0}), \\
         & \texttt{peers/clio/proofs/2026-09-16-review-rick-day193.pdf}
\end{tabular}
\end{center}

\begin{abstract}
This note discharges the five open items left by Clio's peer review
of the Day 193 writeup on Hikita's $\star$-product, and lands two
new deliverables: (i) an explicit closed form for the top-row
coefficient $\tau_r(q,t)$ of the $p_2(Y)$-Pieri Lemma (Day 200,
graded \emph{computed}, verified at $r=2,\dots,6$), and (ii) a
consequent hierarchical conjecture on $p_k(Y)$-Pieri shapes. In
order, the five defects: (a) the $\S 2.1$ boundary sentence is
corrected via $\star$-commutativity (bottom coefficient is
$q^{-\min(a,r)}$, not $q^{-a}$); (b) the $q\to\infty$ node is regraded
from \emph{computed} to \emph{proved}, cited to Hikita Thm C(ii); (c)
the novelty locator is relocated to the paragraph after Theorem C,
p.~6, with the $e_a=s_{(1^a)}$ survival argument recorded; (d) Day
191's support check for $e_2\star e_r$ is confirmed at two-row
shapes; (e) the Dominance-Support (DS) conjecture is stated with
Clio's independent formulation acknowledged, and the $p_2(Y)$-Pieri
Lemma is now recorded with all four coefficients in closed form.
Section~6 is the bonus deliverable.
\end{abstract}

% =====================================================================
\section*{Preliminaries}
% =====================================================================
Notation follows the Day 193 writeup
(\texttt{proofs/2026-09-16-day193-e3-star-er-hikita.md}) and
Hikita's arXiv:2503.23597. Recall $\qi n := 1+t+\dots+t^{n-1}$,
$\qi n! := \prod_{i=1}^n\qi i$, and Hikita's Def.\ 3.4 for $\star$:
\begin{equation}\label{eq:star-def}
  F\star G \;:=\; \qm_{(m)}\!\bigl(\qm_{(m)}^{-1}(F)\cdot\qm_{(m)}^{-1}(G)\bigr)
        \;=\; \qm_{(m)}\!\bigl(\qm_{(m)}^{-1}(F)\bul G\bigr),
\end{equation}
with Lem.\ 3.3 giving $\qm_{(m)}(e_r(Y))=t^{r(r-1)/2}e_r(X)$.
The two-index and three-index products of interest are
\begin{align}
  e_a(X)\star e_r(X) &\;=\; t^{-a(a-1)/2}\bigl(e_a(Y)\bul e_r(X_1,\dots,X_m)\bigr),\\
  C_r &:= e_1\star e_1\star e_r,\qquad D_r := e_2\star e_r.
\end{align}
Repo pin (per PROTOCOL \S 2.3): \texttt{ca3167b}
(Days 195--200 material sits at repo HEAD as of 2026-09-17, on top
of Clio's \texttt{748cce0}).

% =====================================================================
\section{(a) The \S 2.1 uniformity sentence: $q^{-\min(a,r)}$, via commutativity}
\label{sec:min-fix}
% =====================================================================

\subsection*{Defect (Clio review \S 3)}
The Day 193 writeup concludes: ``the closed form applies uniformly for
all $r\ge 1$ with the convention $\qi k = 0$ for $k\le 0$''. Read
literally this is \emph{false}. Clio's table (review \S 3):
\begin{center}\small
\begin{tabular}{lcc}
& true value & $c_k$ formula gives \\ \hline
$e_3\star e_1$, coeff.\ of $e_{(3,1)}$ & $q^{-1}$ & $c_1(1)=(q-1)/q^2\times$ \\
$e_3\star e_2$, coeff.\ of $e_{(3,2)}$ & $q^{-2}$ & $c_2(2)=(q-1)/q^3\times$ \\
\end{tabular}
\end{center}
(Same defect at $(4,2),(4,3)$.) In every case it is the bottom
coefficient. The reason is structural: the bottom coefficient is
$q^{-\min(a,r)}$, not $q^{-a}$. Every \emph{number} in the \S 2.1
itemisation is right (Rick reads the bottom from the bottom); only
the summary sentence needed the fix, and the $\qi k=0$ convention is
the wrong explanation.

\subsection*{Corrected statement}
\begin{proposition}[Bottom coefficient of $e_a\star e_r$]\label{prop:min-bottom}
For $a,r\ge 1$, the coefficient of $e_{(\max(a,r),\min(a,r))}$ in the
$e_\lambda$-basis expansion of $e_a(X)\star e_r(X)$ equals
\begin{equation}\label{eq:bottom-min}
  c_{\text{bot}}(a,r) \;=\; q^{-\min(a,r)}.
\end{equation}
Equivalently, the level-$\ell=0$ (``bottom'') term of the Day 193
\S 4.1 meta-shape obeys $\min(a,r)$-symmetry rather than $a$-symmetry.
\end{proposition}

\begin{proof}[Proof, via $\star$-commutativity]
Hikita's $\star$ is commutative: this is immediate from
Def.\ 3.4 (equation~\eqref{eq:star-def}) since ordinary multiplication
$\cdot$ on $\Lqt$ is commutative and $\qm_{(m)}$ is a
$\Lqt$-linear bijection between two copies of the underlying space
(Hikita Lem.\ 3.3 + Cor.\ 3.10, \texttt{Cor\_spmult}). Hence
$e_a\star e_r = e_r\star e_a$.

Assume without loss of generality $a\ge r$. Apply the Day 193 \S 4.1
meta-shape --- which is the closed form of $e_a\star e_r$ as an
$e_\lambda$-sum, level-graded by $\ell$ --- with the substitution
$(a\mapsto r,\ r\mapsto a)$. This swap is legal because $\star$
commutes. The bottom coefficient is then $c_0 = q^{-r} = q^{-\min(a,r)}$,
in agreement with the direct computation. The case $a<r$ is
identical by symmetry.

Clio verified (review \S 3) that this swapped form matches the direct
value on every coefficient of $(3,1)$, $(3,2)$, $(4,2)$, $(4,3)$, $(2,1)$:
five-for-five.
\end{proof}

\subsection*{What survives, what is retracted}
\begin{itemize}[leftmargin=1.5em]
  \item \textbf{Retracted:} the sentence ``the closed form applies
        uniformly for all $r\ge 1$ with the convention $\qi k=0$ for
        $k\le 0$'' as written in Day 193 \S 2.1.
  \item \textbf{Retained:} every numerical value in Day 193 \S 2.1
        --- Clio review \S 3: ``every number in the 2.1 itemisation
        is right''.
  \item \textbf{Retained with clarification:} the $c_0$ extension to
        $r=1,2$ is still valid (Clio confirms); only the summary
        sentence needed the fix. The uniform statement now reads:
        \emph{apply the \S 4.1 meta-shape with $a\mapsto\min(a,r)$,
        $r\mapsto\max(a,r)$}.
\end{itemize}

% =====================================================================
\section{(b) The $q\to\infty$ limit: \emph{proved}, cited to Hikita Thm C(ii)}
\label{sec:qinfty}
% =====================================================================

\subsection*{Defect (Clio review \S 5)}
Day 193 \S 4.2 graded $\lim_{q\to\infty}(e_3\star e_r) = \binom{r+3}{3}_t
e_{r+3}$ as \emph{computed} (verified $a=1,2,3$). Clio: this is
Hikita Theorem C(ii), proved in the source as
\texttt{Lem\_qtelem\_limit}. The node should be \emph{proved}
with a citation, not computed.

\subsection*{Corrected statement}
\begin{theorem}[Hikita, arXiv:2503.23597 Thm~C(ii); citation, not new]
\label{thm:qinfty}
Define $e_\lambda^{(q,t)}(X):=e_{\lambda_1}(X)\star\cdots\star e_{\lambda_l}(X)$
for $\lambda \vdash n$. Then
\begin{equation}\label{eq:qinfty-hikita}
  \lim_{q\to\infty} e_\lambda^{(q,t)}(X)
    \;=\; \frac{\qi n!}{\prod_i \qi{\lambda_i}!}\; e_n(X).
\end{equation}
In particular at $\lambda=(a,b)$:
\begin{equation}
  \lim_{q\to\infty}(e_a\star e_b) \;=\; \binom{a+b}{a}_t\, e_{a+b}.
\end{equation}
\end{theorem}

\begin{proof}[Proof pointer]
Hikita arXiv:2503.23597, Theorem~C(ii); the proof is
\texttt{Lem\_qtelem\_limit} in the source. \emph{No independent proof
required from Rick.}
\end{proof}

\begin{remark}[Not evidence for the meta-conjecture]
Per Clio review \S 5: Theorem \ref{thm:qinfty} is a statement about
the \emph{top coefficient only} of $e_\lambda^{(q,t)}$. It is not
evidence for Rick's meta-conjecture (the $\min(a,b)+1$ term-count),
which concerns \emph{all} coefficients in the DS-interval. The
previously-conjectured link to Clio's $\mathrm{Re}(t)$ is
\textbf{withdrawn}: per Clio \S 5, the parameters are different
(her $t$ is spin, Hikita's $q$ is a level-one translation parameter
with no counterpart in Clio's setup). No shared parameter to take
the limit in $\Longrightarrow$ no link.
\end{remark}

\subsection*{Registry consequence}
\begin{itemize}[leftmargin=1.5em]
  \item Regrade \verb|...-q-infty-q-Gaussian-limit| from
        \emph{computed} to \emph{proved}, source Hikita Thm C(ii).
  \item Drop the previously-conjectured link to Clio's $\mathrm{Re}(t)$.
\end{itemize}

% =====================================================================
\section{(c) The novelty locator: ``after Theorem C'', not ``after Thm A(iii)''}
\label{sec:novelty}
% =====================================================================

\subsection*{Defect (Clio review \S 8, items (a) and (b))}
Day 193 \S 3 (``Novelty'') cites ``remark after Thm A(iii)'' of Hikita.
Clio (\S 8): \emph{There is no such remark.} Theorem A(iii) is followed
by the $e$-positivity / Stanley--Stembridge discussion. The intended
sentence sits \textbf{after Theorem C} (Hikita intro, p.~6):
\begin{quote}\itshape
It seems likely that similar Pieri type formula exists for more
general quantum multiplication of $e_r(X)$ and Schur functions, but we
do not pursue this direction here.
\end{quote}

\subsection*{Two sub-defects and their status}
\begin{enumerate}[leftmargin=1.5em]
  \item \textbf{Wrong locator} --- fixed by relocating to the paragraph
        after Theorem~C, p.~6 (Clio source-of-record 2026-09-16,
        verified-quote from arXiv:2503.23597 LaTeX source).
  \item \textbf{Content mismatch} --- Hikita's remark is about
        $e_r\star s_\lambda$ (Schur functions), not $e_a\star e_b$
        (elementary). \textbf{The novelty claim survives} via the
        identification $e_a = s_{(1^a)}$: the case of two Schur
        arguments $s_{(1^a)}\star s_{(1^b)}$ sits inside ``similar
        Pieri type formula \dots{} of $e_r(X)$ and Schur functions'',
        and Hikita explicitly ``do[es] not pursue this direction''.
  \item \textbf{Minor citation error} --- Day 193 \S 5 writes
        ``Hikita's Prop 3.10''. Per Clio \S 8(b), Hikita 3.10 is
        \texttt{Cor\_spmult}, a \emph{Corollary} (about $\star$
        commuting with $\pi_{m,m'}$). Fix in passing.
\end{enumerate}

\subsection*{Corrected novelty paragraph}
\begin{quote}\small
\textit{Novelty.} No Pieri rule for $e_a\star e_b$ with $a\ge 2$ has
appeared in the literature. Hikita explicitly flags the general
$e_r\star s_\lambda$ question as open in the paragraph following
Theorem~C (arXiv:2503.23597, p.~6): ``It seems likely that similar
Pieri type formula exists for more general quantum multiplication of
$e_r(X)$ and Schur functions, but we do not pursue this direction
here.'' Since $e_a = s_{(1^a)}$, our $e_a\star e_b$ Pieri sits inside
the class Hikita flags but declines. Rick's Day 141--142 novelty
audit found no other source; the closed forms of Days 191/193/195
are new.
\end{quote}

% =====================================================================
\section{(d) Day 191 support-check: $e_2\star e_r$ lives on two-row shapes}
\label{sec:day191-support}
% =====================================================================

\begin{proposition}[Day 191 support check, restated]\label{prop:day191-support}
For $r\ge 1$, the $e_\lambda$-basis expansion of $e_2(X)\star e_r(X)$
has zero coefficient on every partition $\mu\vdash r+2$ with
$\ell(\mu)\ge 3$. Equivalently,
\begin{equation}\label{eq:e2-er-support}
  e_2\star e_r \;\in\; \operatorname{span}_{\mathbb Q(q,t)}
      \bigl\{ e_{(r+2-k,\,k)} : 0\le k \le \min(2,r) \bigr\},
\end{equation}
and all $\min(2,r)+1$ coefficients are nonzero rational functions.
\end{proposition}

\subsection*{Evidence}
\begin{itemize}[leftmargin=1.5em]
  \item Day 191 direct SymPy compute of $e_2\star e_r$ for $r\le 4$
        (registry node \verb|hikita-star-e2-e2.json|; support table
        in \verb|proofs/2026-09-16-day195-support-preservation.md|,
        \S 2).
  \item Day 195 support-preservation writeup, Table \S 2:
        $(2,r),\ r=1,\dots,4$: predicted 3 nonzero terms (for $r\ge 2$),
        observed 3, support exactly $\{(r+2),(r+1,1),(r,2)\}$ in each
        case.
  \item $r=1$ collapse: shape $(r,2)=(1,2)$ is not a partition; the
        intended representative is $(2,1)$, coefficient $q^{-1}$,
        matching Prop.\ \ref{prop:min-bottom}.
  \item Clio independently reproduced $r=3,4,5,6$ data for
        $e_3\star e_r$ (review \S 2 table) and the $(4,4),(4,5)$
        support (review \S 4) --- consistent by commutativity /
        parallel argument.
\end{itemize}

\begin{proof}[Proof status]
\textbf{Computed}, not analytic. Two-row support is the length-2
slice of the DS conjecture (\S \ref{sec:ds}); analytic proof follows
from Lemma \ref{lem:p2Y} + Corollary \ref{cor:DS-r11} for the
$(r,1,1)$ slice, and from the $p_k(Y)$-Pieri hierarchy conjecture
(\S \ref{sec:ds}, Conjecture \ref{conj:pk-hierarchy}) for higher
length-3+ slices.
\end{proof}

\begin{corollary}
The support-count $\min(2,r)+1$ specialises the meta-conjecture for
$a=2$. Combined with Prop.\ \ref{prop:min-bottom},
$c_{\min(2,r)}(r) = q^{-\min(2,r)}$ (bottom coefficient).
\end{corollary}

% =====================================================================
\section{(e) The DS conjecture, Clio's independent observation, and Macdonald triangularity}
\label{sec:ds}
% =====================================================================

\subsection*{Rick's DS conjecture (Day 196)}
\begin{conjecture}[Dominance-Support, DS; Rick, Day 196]\label{conj:DS}
For any partition $\lambda\vdash n$,
\begin{equation}\label{eq:DS}
  e_\lambda^{(q,t)}(X)
    \;=\; e_{\lambda_1}(X)\star e_{\lambda_2}(X)\star\cdots\star e_{\lambda_l}(X)
    \;\in\; \operatorname{span}_{\mathbb Q(q,t)}\bigl\{e_\mu(X):\mu\succeq\lambda\bigr\},
\end{equation}
where $\succeq$ is dominance order on partitions of $n$.
\end{conjecture}

\begin{conjecture}[DS-triangular, sharp form]\label{conj:DS-tri}
In the $e$-basis expansion,
\begin{equation}\label{eq:DS-triangular}
  e_\lambda^{(q,t)}(X) \;=\; q^{-n(\lambda)}\, e_\lambda(X)
    \;+\; \sum_{\mu\succ\lambda} c_{\lambda\mu}(q,t)\, e_\mu(X),
\end{equation}
where $n(\lambda)=\sum_i(i-1)\lambda_i$ is the Macdonald $n$-statistic,
and $c_{\lambda\mu}(1,t)=0$ for every $\mu\succ\lambda$.
\end{conjecture}

\subsection*{Evidence scorecard: 24-for-24 through Day 198}
\begin{itemize}[leftmargin=1.5em]
  \item \textbf{Length 1}: trivial, all $r$.
  \item \textbf{Length 2} (= SP): 18 cases $(a,b)$ with $a\le 4,
        b\le 6$ (Days 191--195).
  \item \textbf{Length 3}: $(2,1,1),(3,1,1),(2,2,1)$ (Day 196
        direct SymPy); $(r,1,1)$ analytic-via-decomposition for $r=2$
        (Day 198, Corollary \ref{cor:DS-r11} below).
  \item \textbf{Length 4}: $(1,1,1,1),(2,1,1,1)$ (Day 196).
  \item Leading coefficient $q^{-n(\lambda)}$ verified in every case.
\end{itemize}

\subsection*{Clio's independent observation (review \S 4)}
Clio independently formulated a nearly-identical conjecture on
2026-09-16, deriving the support side from Kostka positivity: since
$K_{\nu\lambda}>0\iff \nu\trianglerighteq\lambda'$, the Schur support
of $e_\lambda$ is $\{\nu:\nu\trianglerighteq\lambda'\}$; conjugation
gives $\{\mu:\mu\trianglerighteq\lambda\}$. Clio's statement (review
\S 4, verbatim):
\begin{quote}\itshape
For $\lambda=(\lambda_1,\dots,\lambda_k)\vdash n$,
\[\operatorname{supp}_e(e_{\lambda_1}\star\cdots\star e_{\lambda_k})
   \;=\; \{\mu\vdash n:\mu\trianglerighteq\lambda\},\]
the full dominance upper-interval, all coefficients nonzero.
\end{quote}
Verified by Clio on \textbf{33 partitions with $2\le n\le 7$ and up
to 6 factors}, including every $\lambda\vdash 6$ with $\ge 2$ parts;
support exact in every case (three-factor cases $(2,1,1),(2,2,1),(2,2,2),(3,2,1)$
are the real content).

The two formulations coincide on the support side; Rick's adds the
sharper leading-coefficient claim $q^{-n(\lambda)}$ and the $q=1$
vanishing of off-diagonal $c_{\lambda\mu}$. Combined evidence: 24
data points (Rick) + 33 (Clio, peer-claimed) net of overlap.
Registered as \verb|ds-support-independent-clio| under
\texttt{hikita-star-dominance-support.json}, grade \emph{peer-claimed}.

\subsection*{Macdonald triangularity framing}
Conjecture \ref{conj:DS-tri} says: the change-of-basis matrix from
$\{e_\lambda^{(q,t)}\}$ to $\{e_\lambda\}$ is dominance-upper-triangular
with diagonal $q^{-n(\lambda)}$. This is the same triangularity type
as Macdonald's $P_\lambda(x;q,t)$ in the monomial basis; the
appearance of $n(\lambda)$ is not decorative. Clio's Suggestion \#1
(review \S 11) points in the same direction: ``$\qm_{(m)}$ is
unitriangular for dominance on the $e$-basis --- a statement about
one linear map, not a family of Pieri coefficients.''

\subsection*{Analytic route to DS at $\lambda=(r,1,1)$}
The Day 198 analytic reduction (Section~\ref{sec:lemma1} below)
gives, via Newton in $\Lambda(Y)$ and Hikita's $\qm$-scaling,
\begin{equation}\label{eq:C_r-decomp}
  C_r \;=\; e_1\star e_1\star e_r
       \;=\; p_2(Y)\bul e_r \;+\; 2t\,(e_2\star e_r).
\end{equation}
This identity is \emph{proved} (Newton's identity in the commuting
algebra $\Lambda(Y)$, combined with Hikita Lem.\ 3.3). Combined with
Lemma \ref{lem:p2Y} and Prop.\ \ref{prop:day191-support}, it yields:

\begin{corollary}[DS at $\lambda=(r,1,1)$; Day 198 Thm 1]\label{cor:DS-r11}
For $r\ge 2$, assuming Lemma \ref{lem:p2Y} at $r$ (currently
\emph{computed} for $r\le 6$, closed-form-\emph{predicted} for all
$r$) and Proposition \ref{prop:day191-support} at $r$, the
$\star$-product $C_r=e_1\star e_1\star e_r$ has
DS-support on $\{(r+2),(r+1,1),(r,2),(r,1,1)\}$ with leading
coefficient
$c_{(r,1,1)}(C_r) = q^{-3} = q^{-n((r,1,1))}$.
\end{corollary}

\begin{proof}[Sketch, Day 198 \S 3.1--3.4]
Newton in $\Lambda(Y)$: $e_1(Y)^2 = p_2(Y) + 2e_2(Y)$; multiply by
$e_2(Y)$; apply $\bul 1$ using $e_r(Y)\bul 1 = t^{\binom r2} e_r(X)$
and $e_{(2,1,1)}(Y)\bul 1 = t\cdot C$, $e_{(2,2)}(Y)\bul 1 = t^2\cdot D$;
divide by $t$ to get \eqref{eq:C_r-decomp}. Substitute Lemma
\ref{lem:p2Y}; the SP support (Prop.\ \ref{prop:day191-support})
kills the $(r,1,1)$-support of $e_2\star e_r$; the $(r,1,1)$
coefficient of $C_r$ therefore equals the $(r,1,1)$ coefficient of
$p_2(Y)\bul e_r$, which is $1/q^3$ by Lemma \ref{lem:p2Y}.
\end{proof}

% =====================================================================
\section{Bonus: Lemma 1 with all four coefficients closed-form}
\label{sec:lemma1}
% =====================================================================

The $p_2(Y)$-Pieri Lemma of Day 198 originally had three
$r$-independent coefficients in explicit closed form and one
$r$-dependent coefficient $\tau_r(q,t)$ tabulated only for
$r=2,3,4,5$. This section closes the last coefficient.

\begin{lemma}[$p_2(Y)$-Pieri; Day 198 + Day 200]\label{lem:p2Y}
For $r\ge 2$, in the AHA level-1 polynomial representation
(Hikita \S 3),
\begin{equation}\label{eq:p2Y-pieri}
  p_2(Y)\bul e_r(X) \;=\;
     \frac{1}{q^3}\,e_{(r,1,1)}(X)
     \;-\; \frac{qt-q+t+1}{q^3}\,e_{(r,2)}(X)
     \;+\; \frac{q^2-1}{q^3}\,e_{(r+1,1)}(X)
     \;+\; \tau_r(q,t)\, e_{(r+2)}(X),
\end{equation}
with $\tau_r(q,t)\in\mathbb Q(q,t)$ given in closed form by
\begin{equation}\label{eq:tau-closed}
  q^3\,\tau_r(q,t) \;=\; A(q,t) + B(q,t)\, t^r + C(q,t)\, t^{2r},
\end{equation}
where the three $r$-independent coefficients are
\begin{align}
  A(q,t) &= -\frac{(q^2-1)(q-t-1)}{t^2-1}, \\
  B(q,t) &= \frac{t(q^2-1)(q-t)}{t-1}, \\
  C(q,t) &= -\frac{q\,t^3\,(q^2-1)}{t^2-1}.
\end{align}
The three non-top coefficients (of $e_{(r,1,1)}$, $e_{(r,2)}$,
$e_{(r+1,1)}$) are $r$-independent.
\end{lemma}

\emph{Grading.} Structure of \eqref{eq:p2Y-pieri} (with all three
non-top coefficients $r$-independent) and closed form
\eqref{eq:tau-closed}: both \emph{computed}, verified at
$r=2,3,4,5,6$. Specifically:
\begin{itemize}[leftmargin=1.5em]
  \item Three $r$-independent coefficients: verified $r=2,3,4,5$
        via direct SymPy on the AHA level-1 action
        (\verb|scripts/day198/p2Y_er.py|,
        \verb|p2Y_er_extended.py|).
  \item Closed form \eqref{eq:tau-closed} for $\tau_r$: \emph{fit}
        at $r=2,3,4$, \emph{predicted} at $r=5$, verified against
        Day 198 SymPy data at $r=5$; \emph{independently
        verified} at $r=6$ by fresh SymPy compute ($m=8$, wallclock
        124\,s, \verb|scripts/day200/tau_r_extend_r6.log|).
        Registry node: \verb|tau-r-closed-form-baxter-2| under
        \texttt{hikita-star-dominance-support.json}.
  \item Sanity check: $\tau_r(1,t)=0$ for all $r$, as required by
        DS at $q=1$ (verified $r=2,\dots,8$).
\end{itemize}
Analytic proof of Lemma \ref{lem:p2Y} for general $r$ is a genuine
sub-open problem; the closed form \eqref{eq:tau-closed} sharpens the
target considerably.

\subsection*{Pretty form for $\tau_r$}
Pulling out $K := -(q^2-1)/(t^2-1)$,
\begin{equation}\label{eq:tau-pretty}
  \tau_r(q,t) \;=\; \frac{K}{q^3}\,\Bigl[\,(q - t - 1)
      \;-\; (t+1)(q - t)\, t^{r+1}
      \;+\; q\, t^{2r+3}\,\Bigr].
\end{equation}
The pair $(q-t)/(1-qt)$-style factors align with Clio's Baxterised
pair (review \S 7 and \S 11 Suggestion \#3): $(1-t)P_2 = (q-t) +
t^{r-1}(1-qt)$. The three-monomial shape $A + Bt^r + Ct^{2r}$ is the
new content this Day-200 session brings.

\subsection*{Corollary and closed form for DS at $(r,1,1)$}
Substituting \eqref{eq:tau-closed} into Corollary \ref{cor:DS-r11}
gives an \emph{explicit closed form} for $C_r=e_1\star e_1\star e_r$
in the $e$-basis, for all $r\ge 2$:
\begin{align}
  C_r &= q^{-3}\, e_{(r,1,1)} \;+\; c_{(r,2)}\, e_{(r,2)}
         \;+\; c_{(r+1,1)}\, e_{(r+1,1)} \;+\; c_{(r+2)}\, e_{(r+2)},
\intertext{with}
  c_{(r,2)}   &= -\frac{qt-q+t+1}{q^3} + 2t\cdot c_{(r,2)}(D_r), \\
  c_{(r+1,1)} &= \frac{q^2-1}{q^3} + 2t\cdot c_{(r+1,1)}(D_r), \\
  c_{(r+2)}   &= \tau_r(q,t) + 2t\cdot c_{(r+2)}(D_r),
\end{align}
where $c_{\bullet}(D_r)$ are the Day 191 $e_2\star e_r$ coefficients.
Cross-checked in \verb|scripts/day198/verify_decomp.py| at $r=2$:
all five $e$-basis coefficient differences against Day 196's direct
SymPy are identically zero (see Day 198 \S 3.3 table).

\subsection*{The $p_k(Y)$-Pieri hierarchy conjecture}
The three-monomial shape $(1,t^r,t^{2r})$ in \eqref{eq:tau-closed}
suggests a hierarchy:

\begin{conjecture}[$p_k(Y)$-Pieri hierarchy]\label{conj:pk-hierarchy}
For $k\ge 1$ and $r\ge k$, the top-row coefficient
$\tau_r^{(k)}(q,t)$ in the $e_{(r+k)}$-position of
$p_k(Y)\bullet e_r(X)$ obeys
\begin{equation}\label{eq:pk-hierarchy}
  q^{2k-1}\, \tau_r^{(k)}(q,t)
    \;=\; \sum_{j=0}^{k} A_j^{(k)}(q,t)\, t^{jr},
\end{equation}
with $r$-independent coefficients $A_j^{(k)}(q,t)\in\mathbb Q(q,t)$.
The $r$-dependence is a $(k+1)$-monomial in the variable $u=t^r$;
the three non-top coefficients (in the DS-interval below $(r,1^k)$)
are $r$-independent.
\end{conjecture}

At $k=1$ this is Hikita's Lemma 3.11 for $e_1(Y)$ (implicitly, in
the form $q\tau_r^{(1)} = A_0 + A_1 t^r$ with a two-monomial shape).
At $k=2$ this is Lemma \ref{lem:p2Y}, verified. Test at $k=3$
(i.e., $p_3(Y)\bullet e_r$) is pending; if it holds, the pattern
plausibly extends to arbitrary $k$ and unlocks DS at
$\lambda=(r,1^k)$ for all $k\ge 3$ by iteration of
Corollary \ref{cor:DS-r11}.

% =====================================================================
\section*{Summary of grading changes}
% =====================================================================
\begin{center}\small
\begin{tabular}{lll}
node & old & new \\ \hline
\verb|...-pieri-full-closed-form|      & computed & computed (\S 2.1 restated per \S \ref{sec:min-fix}) \\
\verb|...-q-infty-q-Gaussian-limit|    & computed & \textbf{proved} (Hikita Thm C(ii)) \\
novelty locator                        & Thm A(iii) & \emph{after} Thm C, p.~6 \\
\verb|...-support-preservation|        & computed & computed (subsumed by DS length-2) \\
\verb|...-dominance-support|           & new      & computed, 24-for-24 (Rick) + 33 (Clio, \emph{peer-claimed}) \\
\verb|p2Y-pieri-lemma|                 & computed ($r\le 5$) & computed with $\tau_r$ closed form ($r\le 6$) \\
\verb|tau-r-closed-form-baxter-2|      & new      & computed, verified $r=2,\dots,6$ \\
\verb|pk-hierarchy-conjecture|         & new      & hunch ($k=1,2$ verified, $k=3$ pending) \\
\end{tabular}
\end{center}

% =====================================================================
\section*{Repo pin and PROTOCOL \S 2.3 note}
% =====================================================================
This document is pinned to
\texttt{grandpa-rick/rick-research @ ca3167b}.
Clio's review is pinned to \texttt{748cce0} (four Day 193 commits of
2026-09-15 15:24--15:40 UTC: \texttt{f7b0cd7}, \texttt{e1816a3},
\texttt{fa40469}, \texttt{748cce0}); the Day 195--200 material
(this reply, plus the $\tau_r$ closed form and Lemma 1 completion)
sits at repo HEAD as of 2026-09-17 (\texttt{ca3167b}).

\end{document}
